\documentclass[11pt]{article}
\usepackage[margin=25mm]{geometry}
\usepackage{fontspec}
\setmainfont{Latin Modern Roman}
\usepackage{amsmath,amssymb,mathtools}
\usepackage{unicode-math}
\setmathfont{Latin Modern Math}
\usepackage{xurl,hyperref}
\hypersetup{colorlinks=true,urlcolor=blue}
\setcounter{secnumdepth}{0}
\setlength{\parindent}{0pt}
\setlength{\parskip}{5pt}
\setlength{\emergencystretch}{3em}
\newcommand{\sref}[2]{\hyperlink{source:#1:#2}{[S#2]}}
\newcommand{\cataloguescope}[1]{\paragraph{Recorded target.}#1}
\newcommand{\R}{\mathbb R}
\newcommand{\C}{\mathbb C}
\newcommand{\N}{\mathbb N}
\newcommand{\Z}{\mathbb Z}
\newcommand{\Q}{\mathbb Q}
\newcommand{\F}{\mathbb F}
\newcommand{\eps}{\varepsilon}
\begin{document}
% UnsolvedMath ID 30003911; source catalogue code OWR-16411-003
\setcounter{section}{30003910}
\section{Horizon Topology Changes and Stable Higher-Dimensional Black Holes}\label{30003911}
{\small UnsolvedMath ID 30003911 · Mathematical Physics}
\subsection{Definitions and mathematical statement}

\paragraph{Shared notation.}$\mathbb N=\{1,2,\ldots\}$; the standard number fields have their usual meanings. All additional conventions are stated locally.


\paragraph{Definitions and conventions.}
The Gregory--Laflamme instability is an instability of extended higher-dimensional black objects associated with long-wavelength perturbations. Horizon topology is the topology of a horizon cross-section; a topology-changing evolution must distinguish apparent and event horizons. Myers--Perry metrics are the higher-dimensional rotating vacuum black-hole family. Stability requires a chosen perturbation class. In the source formula, $M$, $J$ and $D$ denote mass, angular momentum and spacetime dimension, and $O(1)$ denotes a bound in its chosen scaling convention.

\paragraph{Target retained from the canonical source catalogue.}
\paragraph{Catalogue wording retained for traceability.}
Is the Gregory–Laflamme instability the only mechanism by which general relativity changes horizon topology? In spacetime dimension $D>4$, is every stable black hole a Myers–Perry black hole satisfying $J/M^{D-3}=O(1)$?

\paragraph{Source scope and qualification.}
Both questions are retained as source conjectural speculation. The supplied $J/M^{D-3}$ is not replaced by a guessed dimensionless normalization; conventional dimensionless spin uses dimension-dependent powers. The asymptotic and rotation hypotheses need source clarification.
\subsection{Short English statement}
Are there other mechanisms for horizon topology changes, and are all stable higher-dimensional black holes moderately rotating Myers--Perry solutions?
\subsection{Research attempt}
\paragraph{Model, date and investigated scope.}
OpenAI GPT-6 Astra (configured model \texttt{gpt-6-astra}), ultra reasoning; generated 11 October 2026. Derived a conditional scaling obstruction to a scale-independent reading of the quoted spin ratio; neither bundled geometric classification is settled.

\paragraph{Sources and literature check.}
Read the supplied definition, catalogue wording, source qualifications and reference entries for 30003911. The historical source/context and inherited status are retained; no fresh literature or novelty certification is claimed. No external theorem is used without proof in this attempt; explicit model assumptions are distinguished from facts about the source model.

\paragraph{A scaling test of the quoted spin bound.} Keep the source expression \(J/M^{D-3}\) exactly as written. Under a vacuum length rescaling by \(s>0\), the usual asymptotic mass and angular momentum coefficients scale as \(M\mapsto s^{D-3}M\), \(J\mapsto s^{D-2}J\): these are the powers of the leading mass and rotation terms in the asymptotic metric. Consequently
\[
 \frac{J_s}{M_s^{D-3}}=s^{D-2-(D-3)^2}\frac{J}{M^{D-3}}.
\]
For every integer \(D>4\), the exponent is negative: it is \(-1\) at \(D=5\) and decreases thereafter. Thus, conditional on a scale-invariant stable family with nonzero spin, this ratio has no uniform scale-independent upper bound as \(s\downarrow0\). In contrast \(J/M^{(D-2)/(D-3)}\) is invariant under this calculation.
\paragraph{Scope and obstruction.} This is a normalization audit, not a silent replacement of the source bound and not a black-hole stability theorem. A fixed mass unit, compact parameter range or an intended different convention can change the interpretation of \(O(1)\). The extraction explicitly leaves that point unresolved. Even after it is fixed, scaling identifies neither all stable solutions nor a topology-change mechanism. The Gregory--Laflamme exclusivity question and the Myers--Perry classification both remain. The first gap is a precise stability/asymptotic class and a scale-qualified formulation on which either a dynamical argument or a counterexample could be tested.

\paragraph{Verification and adversarial audit.}
Checked the exponent for all D greater than four and avoided asserting existence/stability of a nonzero-spin counterexample family. This is a GPT-6 Astra Ultra self-audit, not independent human verification. Prompt A drafting followed by prompt B scope/dependency/boundary self-audit. Checked the exponent for all D greater than four and avoided asserting existence/stability of a nonzero-spin counterexample family. The remaining gap is: Fix the source scale and stability conventions, then analyze actual stable solutions and topology-changing evolutions. Unresolved status and zero original-target completion are retained.

Coordinator checked both scaling powers and the sign of the resulting exponent for every integer dimension greater than four. The argument is conditional on the stated stable scaling family. This additional check was model-assisted, not independent human verification.

\paragraph{Tools, process and contributions.}
Python 3 and local file-reading/writing tools for manuscript serialization and structural checks; mathematical calculations were symbolic reasoning, with no numerical experiment or theorem prover. Contribution prompt A was followed by prompt B. GPT-6 Astra with ultra reasoning drafted the research and performed its stated self-audit; any additional coordinator contribution is identified in the verification and audit. A GPT-6-based Codex coordinator assembled the manuscript and checked repository integration; its exact variant and reasoning configuration were not exposed. No human mathematical contribution, independent human proof check or formal proof certification is supplied. Complete derivations and the consulted source references are included here; private reasoning transcripts are not published.

\paragraph{Final status and first remaining gap.}
\textbf{UNRESOLVED}. Fix the source scale and stability conventions, then analyze actual stable solutions and topology-changing evolutions.

\section{Completion Estimate}
\noindent\textbf{COMPLETION: 0\%}
The proved calculation is an elementary auxiliary result or a test of a restricted strategy; no new portion of the full source target is established.

\subsection{Sources}
\begin{itemize}
\item[\hypertarget{source:30003911:1}{S1}] ULAM.AI and the UnsolvedMath Contributors, UnsolvedMath, record 30003911 (OWR-16411-003). Catalogue curation under CC BY 4.0; definitions and English exposition added locally.
\url{https://huggingface.co/datasets/ulamai/UnsolvedMath}.
\item[\hypertarget{source:30003911:2}{S2}] Mathematical General Relativity (2018), Oberwolfach Reports, DOI 10.4171/owr/2018/36, pp. 2157--2251. Primary report accessed 2026-10-11; relevant PDF pages 37.
\url{https://doi.org/10.4171/owr/2018/36}.
\item[\hypertarget{source:30003911:3}{S3}] Emparan--Reall higher-dimensional black holes (catalogue-listed context; not independently consulted in this preparation).
\url{https://doi.org/10.1007/978-3-642-41992-8}.
\item[\hypertarget{source:30003911:4}{S4}] Black holes galore in D greater than 4 (catalogue-listed context; not independently consulted in this preparation).
\url{https://doi.org/10.1007/s10714-025-03398-x}.
\end{itemize}
\label{end:30003911}
\end{document}
